\scanpage{4}
\[
\begin{aligned}
(4)\qquad\qquad & \operatorname{Pic}(X)\handwritten{\overset{\sim}{\mathrel{\tikz[baseline=-.3ex,x=1em,y=1em]{\draw[line width=.45pt] (0,0)--(.93,0);\draw[line width=.45pt] (.93,.05) .. controls (1.07,.10) and (1.08,-.14) .. (.96,-.12)--(.93,.05);}}}}P(k)\ .
\end{aligned}
\]
The map $D\handwritten{\longrightarrow}\operatorname{Pic}(X)\handwritten{\simeq}P(k)$ défines a morphism of (pro) \rlap{G}gro\uline{u}ps
\[
\begin{aligned}
(5)\qquad\qquad & D_k\handwritten{\longrightarrow}P\ ,
\end{aligned}
\]
where $D_k$ is the constant group-scheme with value $D$. Assuming that the cokernel of (5) exists in any reasonable sense a a pro-group scheme or a group-scheme, let $Q$ be this kernel. We get, in virtue of (1), an injective homomorphism
\[
\begin{aligned}
(6)\qquad\qquad & \operatorname{Pic}(U)\handwritten{\hookrightarrow}Q(k)\ ,
\end{aligned}
\]
the obstruction for an element of $\unsure{Q|\mathbb{Q}}(k)$ to be in the image being in $H^1(k,D'_k)$, where $D'_k$ is the image of $D_k$ in $P$ (it is again a constant group, with value a quotient $D'$ of $D$). \tikz[baseline=(s.base)]{\node[inner sep=1pt](s){\deleted{If}};\node[anchor=south west,inner sep=1pt](a) at ([xshift=-2pt,yshift=7pt]s.north west){\inserted{\handwritten{If $k$ is sep. closed, or if}}};\draw ([xshift=1pt,yshift=-2pt]s.north west) to[out=70,in=190] ([xshift=1pt,yshift=-2pt]a.south west)--([xshift=1pt,yshift=-2pt]a.south east) to[out=0,in=-90] ([xshift=3pt,yshift=5pt]a.south east);} $D'$ is torsion-free, for example is $D'$ is equal to $D$ i.e. (5) is injective, then this obstruction vanishes, and \deleted{$Q$ is a good candidat\unsure{e}} (6) is an isomorphism. Thus $Q$ looks a good candidate for a ``local Picard scheme''.

To construct the $P_n$'s, assume first that $A$ is a $k$-algebra, i.e. that a lifting of \tikz[baseline=(s.base)]{\node[inner sep=1pt](s){\deleted{\unsure{thh|the}}};\node[anchor=south,inner sep=1pt] at ([xshift=5pt,yshift=1pt]s.north){\inserted{the}};} residue field $k$ exists, and has been choosen. Then the natural candidate for $P_n$ (the only one I could think of, indeed !) would be
\[
\begin{aligned}
(7)\qquad\qquad & P_n=\underline{\operatorname{Pic}}_{X_n/k}\ .
\end{aligned}
\]
We know from Murre that the right had side is indeed representable by a scheme locally of finite type over $k$. There is however the usual trouble that we have indeed a canonical \uline{injective} morphism
\[
\begin{aligned}
(8)\qquad\qquad & \operatorname{Pic}(X_n)\handwritten{\hookrightarrow}P_n(k)\ ,
\end{aligned}
\]
but in general we are not sureit is bijective. However it is if $k$ is separably closed; more specifically, the obstruction for an element of $P_n(k)$ te be in the image of $\operatorname{Pic}(X_n)$ lies in the Brauer group \deleted{$\operatorname{Br}(k)$} \par\noindent\deleted{of $\unsure{Bo|B_o|k}\,\tikz[baseline=-.3ex,x=1em,y=1em]{\draw[line width=.6pt] (.10,.30) .. controls (-.04,.10) and (-.04,-.45) .. (.18,-.63);}$}
% archive sheet number: 3
